\scanpage{80}
\mspage{39}
\[
\tag{102}
\diagram{\begin{tikzcd}[ampersand replacement=\&, column sep=large,row sep=large]
\& \Pi M^*_{\underline G,\underline G'} \arrow[r] \arrow[d] \& \Pi M_{G'} \\
\Pi(M_G\deleted{,\unsure{\Gamma^!_{\underline G,\underline G'}}}) \arrow[r,"\simeq"] \& \Pi M_{G,G'}
\end{tikzcd}}
\]
\margin{ceci est une équivalence de groupoïdes}
qu'on va réécrire comme
\[
\tag{103}
\diagram{\begin{tikzcd}[ampersand replacement=\&, column sep=large,row sep=large]
\Pi M^*_{\underline G,\underline G'} \arrow[r,"\unsure{\psi_{\underline G,\underline G'}|\gamma_{\underline G,\underline G'}}"] \arrow[d,"\varphi_{G,G'}"] \& \Pi M_{G'} \\
\Pi(M_G\deleted{,\Gamma^!_{\underline G,\underline G'}})
\end{tikzcd}}
\]
\margin{« diagramme de générisation pour les groupoïdes fondamentaux, \unsure{pour} le quotient de $\underline G\to\underline G'$ »}
\deleted{pour la suite, il sera plus commode d'écrire $\Pi M_{G,G'}$ au lieu de $\Pi(M_G,\Gamma^!_{\underline G,\underline G'})$}

L'hom des groupoïdes $\varphi_{G,G'}$ a les propriétés de \deleted{\unsure{la}} l'hom sur les $\pi_1$ dans une fibration de base $(M_G\deleted{,\Gamma_{\underline G,\underline G'}\inserted{\text{op. de }\Gamma^!_{G,G'}}})$, \deleted{de} fibre un produit de disques épointés (ou de cercles, si on préfère) --- on peut considérer $\Pi M^*_{\underline G,\underline G'}$ comme une « \uline{extension} » canonique des groupoïdes $\Pi(M_G\deleted{,\Gamma^!_{\underline G,\underline G'}})$ par $\mathbf Z^C$, où $C$, on le rappelle, est l'ens. des \deleted{frontières}\inserted{arêtes} de $\underline G$

\margin{Attention, il faut regarder $\mathbf Z^C$ comme un syst. local sur $(M_G,\Gamma_{\underline G,\underline G'})$, via les op. de $\Gamma^!_{\underline G,\underline G'}$ sur $C$\ldots}
\margin{\uline{NB} $(M_G\deleted{,\Gamma_{\underline G,\underline G'}})$ est un $K(\pi,1)$, et par conséquent \unsure{dans|on a dans} la suite exacte des $\pi_1$ de la fibration \unsure{qui précède} \unsure{incluse, on a bien un $1$ à gauche|des inclusions et bien un $1$ à gauche}.}
\margin{79}
